\section{Matter on the Curved Substrate}

To see why the committed substrate has to carry spin in its sites, it helps to study a neighbour that does not. The three-dimensional honeycomb $\{5,3,4\}$ is the nearest relative of the committed mesh with a three-dimensional cusp, and on it the question of matter is sharp, because its sites refuse spin. Its twelve directions form the icosahedral group, and a permutation representation of a finite group is linear while a spinor is projective, so no sum or subrepresentation of the bare field on those directions will ever be a fermion. This is not an engineering difficulty, it is a theorem.

\begin{theorem}[The spin obstruction]
On $\{5,3,4\}$ the bare directional rule carries no spinor. The twelve directions act by a linear representation of the icosahedral group, and a spinor requires a projective one, so no bound state of the bare field is fermionic.
\end{theorem}

And yet the geometry itself still carries the spin structure, at every scale, even though the sites do not, and a field can be made to live in that structure by four independent routes that close the question. The first is to put the fermion on the form complex of the lattice rather than on the bare sites, where the natural operator squares to the Laplacian and is the Kahler-Dirac fermion of lattice field theory. Run on the actual $\{5,3,4\}$ bulk, this fermion genuinely propagates, its return probability low and falling with size, in the extended rather than the trapped phase, and it localizes only when a deterministic quasiperiodic disorder is switched on, a control that could have failed (depth L2). The second route reads spin off a defect. Transporting a probe around an odd disclination in the director field flips the spinor sign while leaving vectors untouched, the topological minus-one of the double cover, carried by the defect with no change to the rule (depth L2). The third route shows the same sign carried by a delocalized collective mode rather than a point probe, and for every mode, so it is a property of the field and not an artifact of localization (depth L2). The fourth is purely geometric, the holonomy around a loop in this honeycomb is the binary icosahedral group, the double cover itself.

So $\{5,3,4\}$ can host matter after all, through its form complex and its defects, but only by a detour, because its bare directions lack what the committed substrate has for free. That contrast is the point. There is a clean way to have both the spinor directions and the curvature at once, and it is to climb one more dimension to the five-dimensional pentacomb $\{3,4,3,3,4\}$, whose directions already carry the spinor and whose space is curved, where a Kahler-Dirac fermion propagates directly with no detour at all (depth L2). The pentacomb proves that spin and curvature need not be traded against each other, and the committed substrate $\{3,4,3,4\}$ is the four-dimensional member that keeps the spinor sites while handing physics a three-dimensional cusp. The comparison substrate earns its place by showing exactly what the thesis substrate provides that it does not.
